% ============================================================
% 纯答案册·课时15-2.7.1抛物线方程 —— 工具/答案抽册器.py 抽册生成件（勿手改；第三档＝纯答案单独成册）
% 源件（只读）：C:/提示词/工作区/M3-第2章量产0913/成卷/练习件/课时15-2.7.1抛物线方程/main.tex（md5 5b38f5a3326bce7f20acaaf5ee3a0ccc）
%% 口径：题面不重印；逐题＝题号(印面号同正册)＋[答案]＋[详解]，值/详解照源件逐字
%       （钉值门硬断言：册值≡正册件内值≡值台账值tex，逐字全等）。册序＝正册装配序＝印面号 1..16 升序（同序同号）；键序断言基准＝题面库冻结 manifest canonical 键序。
%% 三档导出：整体 main-true／纯题 main-false（在产）｜本册＝纯答案册（本件）
%% 双档：ansbook-true.tex＝详解印本；ansbook-false.tex＝纯值速查本（详解不印）
% 编译：xelatex <壳> ×2，cwd＝本目录；sty＝qp-m3.sty 本地挂载（md5 7c3930362be8a0a2bdf21bbf8ac16573，锁校验）
% ============================================================
\documentclass[fontset=none]{ctexart}
\usepackage{qp-m3}
% —— 印面逐字符号路由补钉（承源件；− 走 TNR、▱ NSC 子块；禁改 sty 保 md5） ——
\xeCJKDeclareCharClass{Default}{"2212}%
\setCJKfamilyfont{nscblk}[Path=C:/提示词/工作区/字替对照-0909/variantF/fonts/,BoldFont=NSC-w600.ttf]{NSC-w500.ttf}%
\xeCJKDeclareSubCJKBlock{nscblk}{"25B1}%
\newcommand{\pxparallelogram}{{\CJKfamily{nscblk}▱}}%
% —— 断行微调（承母版实测档，三零门承重；\emergencystretch 不另设） ——
\xeCJKsetup{CJKglue={\hskip 0.10em plus 0.08em minus 0.05em}}%
% —— 纯答案册全线括线（长详解可跨栏断，承练习件全线括线口径；灰底盒不可跨栏断） ——
\ansblockgrayfalse
% —— 双档开关（false 档＝纯值速查：答案印、详解不印；件面开关，不动 sty 本体） ——
\newif\ifansbookdetail
\ifdefined\ansbookdetail
  \ifnum\ansbookdetail=0 \ansbookdetailfalse\else\ansbookdetailtrue\fi
\else
  \ansbookdetailtrue
\fi
% —— 页脚件名（承源件章名；件名＝<件型>答案册） ——
\renewcommand{\qpzhangming}{第二章\quad 平面解析几何}%
\renewcommand{\qpjianming}{练习件答案册}%

\begin{document}
\keshi{第15课时\quad 抛物线的标准方程}
\par\glueguard{1}\addvspace{2pt}\noindent\biaoqian{【纯答案册】}{\fangsong 题面不重印\quad 印面号与正册同号同序}\par\addvspace{2pt}

\begin{multicols}{2}
\emergencystretch=1em

\begin{ansblock}[2章-练-课时15-01]
% ans:2章-练-课时15-01
\ansitem{1}{不是。当定点 F 在直线 l 上时，动点轨迹是与 l 垂直的直线（过 F）。}
\ifansbookdetail
\ansnote{详解}{取\(l\)为\(x\)轴、\(F\)为原点。动点\(P(x,y)\)到\(F\)的距离为\(\sqrt{x^{2}+y^{2}}\)，到\(l\)的距离为\(|y|\)。\par\noindent “到\(F\)与到\(l\)距离相等”：\(\sqrt{x^{2}+y^{2}}=|y|\)，平方得\(x^{2}+y^{2}=y^{2}\)，即\(x=0\)。轨迹为整条\(y\)轴：过\(F\)且垂直于\(l\)的直线，不是抛物线（\(p=0\)时抛物线无意义）。验算：点\((0,5)\)到\(F\)距离5、到\(x\)轴距离5，相等合。}
\fi
\end{ansblock}

\begin{ansblock}[2章-练-课时15-02]
% ans:2章-练-课时15-02
\ansitem{2}{y²=3x 或 y²=−3x 或 x²=3y 或 x²=−3y。}
\ifansbookdetail
\ansnote{详解}{焦准距即\(p\)：\(p=3/2\)，\(2p=3\)。开口未指定，四种标准方程并立：\(y^{2}=3x\)、\(y^{2}=-3x\)、\(x^{2}=3y\)、\(x^{2}=-3y\)。}
\fi
\end{ansblock}

\begin{ansblock}[2章-练-课时15-03]
% ans:2章-练-课时15-03
\ansitem{3}{x²=−12y}
\ifansbookdetail
\ansnote{详解}{\(F(0,-3)\)不在直线\(y=3\)上，由定义轨迹为抛物线：焦点\(F(0,-3)\)，准线\(y=3\)。开口向下，设\(x^{2}=-2py\)（\(p>0\)）：\(p/2=3\)，\(p=6\)，方程\(x^{2}=-12y\)。验算：顶点\((0,0)\)到\(F\)距3、到准线距3合。}
\fi
\end{ansblock}

\begin{ansblock}[2章-练-课时15-04]
% ans:2章-练-课时15-04
\ansitem{4}{y²=−16x}
\ifansbookdetail
\ansnote{详解}{焦点在\(x\)轴，设\(y^{2}=mx\)（\(m\neq 0\)）。代入\((-4,-8)\)：\(64=m\times (-4)\)，\(m=-16\)，方程\(y^{2}=-16x\)（开口向左，与横坐标为负相符）。}
\fi
\end{ansblock}

\begin{ansblock}[2章-练-课时15-05]
% ans:2章-练-课时15-05
\ansitem{5}{C}
\ifansbookdetail
\ansnote{详解}{左端为\(P(x,y)\)到定点\(F(0,p/2)\)的距离；右端\(|y+p/2|\)为\(P\)到定直线\(l\)：\(y=-p/2\)的距离。\(F(0,p/2)\)不在\(l\)上（\(p>0\)），由抛物线定义，轨迹为抛物线，故选C。}
\fi
\end{ansblock}

\begin{ansblock}[2章-练-课时15-06]
% ans:2章-练-课时15-06
\ansitem{6}{3}
\ifansbookdetail
\ansnote{详解}{由抛物线定义，抛物线上点到焦点的距离与到准线的距离相等，故\(M\)到准线的距离为\(|MF|=3\)。}
\fi
\end{ansblock}

\begin{ansblock}[2章-练-课时15-07]
% ans:2章-练-课时15-07
\ansitem{7}{(1) y²=8x；(2) y²=6x。}
\ifansbookdetail
\ansnote{详解}{(1)焦点\((2,0)\)在\(x\)正半轴：\(p/2=2\)，\(p=4\)，\(2p=8\)，方程\(y^{2}=8x\)。(2)准线\(x=-3/2\)得焦点在\(x\)正半轴：\(p/2=3/2\)，\(p=3\)，\(2p=6\)，方程\(y^{2}=6x\)。验算：(1)准线\(x=-2\)合；(2)焦点\((3/2,0)\)合。}
\fi
\end{ansblock}

\begin{ansblock}[2章-练-课时15-08]
% ans:2章-练-课时15-08
\ansitem{8}{y²=24x}
\ifansbookdetail
\ansnote{详解}{焦点在\(x\)正半轴，开口向右，设\(y^{2}=2px\)（\(p>0\)），准线\(x=-p/2\)。准线与\(y\)轴距离\(|-p/2|=6\)得\(p=12\)，\(2p=24\)，方程\(y^{2}=24x\)。验算：焦点\((6,0)\)、准线\(x=-6\)，与\(y\)轴距离均为6合。}
\fi
\end{ansblock}

\begin{ansblock}[2章-练-课时15-09]
% ans:2章-练-课时15-09
\ansitem{9}{y²=8x}
\ifansbookdetail
\ansnote{详解}{设\(y^{2}=2px\)（\(p>0\)）。代入\((2,-4)\)：\(16=4p\)，\(p=4\)，\(2p=8\)，方程\(y^{2}=8x\)。验算：焦点\((2,0)\)、准线\(x=-2\)；\(M(2,-4)\)到焦点距离4、到准线距离4合。}
\fi
\end{ansblock}

\begin{ansblock}[2章-练-课时15-10]
% ans:2章-练-课时15-10
\ansitem{10}{M(6,6√2) 或 M(6,−6√2)。}
\ifansbookdetail
\ansnote{详解}{\(y^{2}=12x\)：\(2p=12\)，\(p=6\)，准线\(x=-3\)。由定义\(|MF|=x_M+3=9\)，得\(x_M=6\)；代入得\(y^{2}=72\)，\(y=\pm 6\sqrt{2}\)。故\(M(6,6\sqrt{2})\)或\(M(6,-6\sqrt{2})\)（两点均在抛物线上，关于\(x\)轴对称）。验算：\(\sqrt{(6-3)^{2}+72}=\sqrt{81}=9\)合。}
\fi
\end{ansblock}

\begin{ansblock}[2章-练-课时15-11]
% ans:2章-练-课时15-11
\ansitem{11}{y²=8x 或 x²=−y}
\ifansbookdetail
\ansnote{详解}{\(A(2,-4)\)在第四象限，开口向右或向下。开口向右：设\(y^{2}=2px\)（\(p>0\)），\((-4)^{2}=2p\times 2\)得\(p=4\)，方程\(y^{2}=8x\)；开口向下：设\(x^{2}=-2py\)（\(p>0\)），\(4=-2p\times (-4)\)得\(p=1/2\)，方程\(x^{2}=-y\)。综上\(y^{2}=8x\)或\(x^{2}=-y\)。验算：\(4^{2}=8\times 2\)合；\(2^{2}=-(-4)\)合。}
\fi
\end{ansblock}

\begin{ansblock}[2章-练-课时15-12]
% ans:2章-练-课时15-12
\ansitem{12}{x=0；1/2。}
\ifansbookdetail
\ansnote{详解}{设\(P(x,y)\)：\(\sqrt{x^{2}+(y-1)^{2}}+|y+1|=3\)。①\(y>2\)或\(y<-4\)（\(|y+1|>3\)）：左边\(>3\)，无解；②\(-1\leqslant y\leqslant 2\)：\(\sqrt{x^{2}+(y-1)^{2}}=2-y\)，平方得\(x^{2}=-2y+3\)（\(-1\leqslant y\leqslant 3/2\)），对称轴\(x=0\)，此时\(|PF|=2-y\geqslant 1/2\)；\par\noindent ③\(-4\leqslant y<-1\)：\(\sqrt{x^{2}+(y-1)^{2}}=y+4\)，平方得\(x^{2}=10y+15\)（\(-3/2\leqslant y<-1\)），对称轴\(x=0\)，此时\(|PF|=y+4\geqslant 5/2\)。综上，对称轴方程\(x=0\)，\(|PF|_{\min}=1/2\)。验算：②段平方展开\(x^{2}+(y-1)^{2}=(2-y)^{2}\)得\(x^{2}=-2y+3\)合；②段与③段区间衔接经复核无缝合。}
\fi
\end{ansblock}

\begin{ansblock}[2章-练-课时15-13]
% ans:2章-练-课时15-13
\ansitem{13}{B}
\ifansbookdetail
\ansnote{详解}{以顶端\(O\)为原点建系，设\(x^{2}=-2py\)（\(p>0\)）。由对称性\(D(15,h)\)（\(h<0\)），\(B(30,h-150)\)。联立：\(225=-2ph\)，\(900=-2p(h-150)\)。两式相除：\(4=(h-150)/h\)，得\(h=-50\)；代入\(225=-2p\times (-50)\)得\(p=2.25\)。顶端\(O\)到\(AB\)的距离\(=|h|+150=200\)(m)，故选B。验算：\(p=2.25\)，\(2p=4.5\)，\(900=4.5\times 200\)合。}
\fi
\end{ansblock}

\begin{ansblock}[2章-练-课时15-14]
% ans:2章-练-课时15-14
\ansitem{14}{y²=16x。}
\ifansbookdetail
\ansnote{详解}{设\(M(x,y)\)。“\(|MF|\)比到\(l\)的距离小2”即\(|x+6|-|MF|=2\)。当\(x\geqslant -6\)时：\(x+6-\sqrt{(x-4)^{2}+y^{2}}=2\)，即\(\sqrt{(x-4)^{2}+y^{2}}=x+4\)（隐含\(x\geqslant -4\)，自然满足），平方得\((x-4)^{2}+y^{2}=(x+4)^{2}\)，化简为\(y^{2}=16x\)（此时\(x\geqslant 0\geqslant -4\)，与所设一致）；\par\noindent 当\(x<-6\)时：\(-x-6-\sqrt{(x-4)^{2}+y^{2}}=2\)，即\(\sqrt{(x-4)^{2}+y^{2}}=-x-4\)，平方同样得\(y^{2}=16x\)，但该抛物线上\(x\geqslant 0\)，与\(x<-6\)矛盾，此段无轨迹。故轨迹方程\(y^{2}=16x\)（顶点在原点、开口向右的抛物线，焦点\(F(4,0)\)、准线\(x=-4\)）。验算：\(y^{2}=16x\)上任一点\(|MF|=x+4\)，到直线\(x=-6\)的距离为\(x+6\)，差恰为2合。\par\noindent 辨析：本题须防“大2”误读——若误作“\(|MF|\)比到直线\(x=-6\)的距离大2”，则\(|MF|=|x+8|\)，平方化简为\(y^{2}=24(x+2)\)，为平移型抛物线（顶点\((-2,0)\)），与标准形不符；按原题“小2”，平移准线得\(|MF|\)等于到直线\(x=-4\)的距离，才落到标准抛物线\(y^{2}=16x\)。}
\fi
\end{ansblock}

\begin{ansblock}[2章-练-课时15-15]
% ans:2章-练-课时15-15
\ansitem{15}{(1) y²=4x；(2) 存在，此时 Q(9,6) 或 Q(4,−4)。}
\ifansbookdetail
\ansnote{详解}{(1)\(16=8p\)得\(p=2\)，\(C\)：\(y^{2}=4x\)。\par\noindent (2)\(F(1,0)\)。设\(N(t,2t+6)\)、\(Q(x_0,y_0)\)。四边形\(PQFN\)为平行四边形\(\Leftrightarrow\overrightarrow{PQ}=\overrightarrow{NF}\Leftrightarrow Q-P=F-N\)：\((4-t,-2t-2)=(x_0-1,y_0)\)，得\(x_0=5-t\)，\(y_0=-2t-2\)。\(Q\)在\(C\)上：\((-2t-2)^{2}=4(5-t)\)，即\(4t^{2}+8t+4=20-4t\)，\(t^{2}+3t-4=0\)，\(t=1\)或\(t=-4\)。\(t=1\)：\(N(1,8)\)，\(Q(4,-4)\)；\(t=-4\)：\(N(-4,-2)\)，\(Q(9,6)\)。\par\noindent 非退化核验（\(P\)、\(Q\)、\(F\)、\(N\)不共线，\(N\neq F\)）：\(t=1\)时\(\overrightarrow{PQ}=(0,-8)=\overrightarrow{NF}\)合；\(\overrightarrow{QF}=(-3,4)\)，\(\overrightarrow{PN}=\overrightarrow{P}-\overrightarrow{N}=(3,-4)\)，平行且相等合；四点不共线合，\(N(1,8)\neq F(1,0)\)合；\par\noindent \(t=-4\)时\(\overrightarrow{PQ}=(5,2)\)，\(\overrightarrow{NF}=\overrightarrow{F}-\overrightarrow{N}=(5,2)\)合；\(\overrightarrow{QF}=(-8,-6)\)，\(\overrightarrow{PN}=(8,6)\)合；斜率\(k_{PQ}=2/5\neq k_{QF}=3/4\)，四点不共线合，\(N\neq F\)合。两解均成立：存在\(N\)使\(PQFN\)为平行四边形，\(Q(9,6)\)或\(Q(4,-4)\)。验算：\(t^{2}+3t-4=(t-1)(t+4)\)合；\(Q(9,6)\)：\(36=4\times 9\)合；\(Q(4,-4)\)：\(16=16\)合。}
\fi
\end{ansblock}

\begin{ansblock}[2章-练-课时15-16]
% ans:2章-练-课时15-16
\ansitem{16}{不能安全通过隧道。}
\ifansbookdetail
\ansnote{详解}{以抛物线顶点为原点、对称轴为\(y\)轴建系。由图：拱高3m、半宽3m，抛物线过\(A(3,-3)\)；矩形部分高2m、总高（顶点至地面）5m、地面\(y=-5\)。设\(x^{2}=-2py\)（\(p>0\)）：\(9=-2p\times (-3)\)，\(2p=3\)，方程\(x^{2}=-3y\)。\par\noindent 车居中：\(x=\pm 1.5\)处拱高\(y=-(1/3)\times 2.25=-0.75\)，净空\(=5-0.75=4.25\)(m)。\(4.25<4.5\)（车与集装箱总高），即使集装箱处于隧道正中，总高仍高于该处拱高，故不能安全通过。验算：\((\pm 1.5)^{2}=-3y\)得\(y=-0.75\)合；\(2p=3\)时准线\(y=3/4\)合。}
\fi
\end{ansblock}

% ---- 尾块（\tailfill[30mm] 定高参——钉档 §三.3：无参在平衡栏呈「内容尾随框」，贴栏底须定高参；实测无参致末页平衡 Overfull\vbox 12.99pt，定高 30mm 三零） ----
\tailfill[30mm]

\end{multicols}
\end{document}
